\documentclass[10pt]{article}
\usepackage{amsmath,amssymb,geometry}
\geometry{margin=1in}
\title{Two maximum antichains disprove conjecture 00000008544}
\author{Local AI-assisted solution draft}
\date{3 October 2026}
% Compact consistent title for a short mathematical note.
\makeatletter
\renewcommand{\maketitle}{\begin{center}
{\Large\bfseries\@title\par}\vspace{0.6em}
{\small\@author\par\@date}
\end{center}\vspace{0.5em}}
\makeatother
\setlength{\emergencystretch}{3em}
\begin{document}
\maketitle
\paragraph{Claim being disproved.}
The conjecture claims that a Cartesian product of total-order chains has a unique maximum antichain. The product of three two-element chains is a counterexample.
\paragraph{Construction.}
Let $P=\{0,1\}^3$ with the coordinatewise order, equivalently the Boolean lattice of subsets of $\{1,2,3\}$. Define
\[
 \mathcal A=\{100,010,001\},\qquad
 \mathcal B=\{110,101,011\}.
\]
Distinct members of either family are incomparable: they have equal Hamming weight, and a coordinatewise inequality at equal weight forces equality. Both families have size three and they are distinct.
\paragraph{Maximum size.}
For completeness, partition all eight elements into the following three chains:
\[
 000<001<011<111,\qquad 010<110,\qquad 100<101.
\]
Every antichain meets each chain at most once, so it has at most three elements. Consequently both $\mathcal A$ and $\mathcal B$ are maximum antichains. They are the two middle rank layers, demonstrating why a uniqueness claim fails in this example. The Sperner number is $3$; its precise expression elsewhere in the conjecture need not be settled to refute uniqueness.
\paragraph{Formal verification and semantic correspondence.}
The accompanying \texttt{Main.lean} uses Lean 4.19.0 \texttt{Std}. The eight elements are indexed by \texttt{Fin 8}; \texttt{cube} and \texttt{unCube} are explicitly checked inverse maps with $(\texttt{Fin 2})^3$. The order \texttt{Below} is exactly the coordinatewise order under this identification, and totality of the factors and the resulting partial-order laws are checked.

Subsets are represented by eight-bit masks \texttt{Fin 256}. Membership is defined by the corresponding bit; \texttt{every\_subset\_encoded} proves that every Boolean membership function on the eight elements has such an encoding, so the subset enumeration is exhaustive. \texttt{Antichain} uses actual membership and comparability, and \texttt{cardinality} counts actual members. The theorem \texttt{sharp\_upper\_bound} checks all $256$ subsets and proves that every antichain has at most three elements. This is an independent finite verification of the chain-partition argument above.

The masks $22$ and $104$ represent $\mathcal A$ and $\mathcal B$ respectively; their member descriptions, sizes, maximum status, and inequality are all checked. The final theorem \texttt{conjecture8544\_false} directly negates uniqueness of a maximum antichain on this chain product. It does not merely exhibit incomparable elements without checking maximal size.

All finite checks use kernel reduction via \texttt{decide}. The final disproof has only the standard axiom \texttt{propext}; the general encoding lemma additionally uses standard \texttt{Quot.sound} through function extensionality. There are no placeholders, \texttt{native\_decide}, or additional axioms.
\paragraph{Reproduction.}
Run \texttt{lean Main.lean} with Lean 4.19.0.
\end{document}
